% SOURCE_AUTHORITY: sga5_fr_workpass.tex, lines 8695--8705 (LF-normalized unit).
% SOURCE_SHA256: 791F4EFFC5E02832D5D77ED03518C8156D6F07E4C8238B03545DB93D883FBB28
\subsubsection*{1.4.7}
Seja agora \(L\) uma extensão de tipo finito de \(\mathbb Q_\ell\), e denotemos por \(A\) a clausura integral de \(\mathbb Z_\ell\) em \(L\), de modo que \(A\) é uma \(\mathbb Z_\ell\)-álgebra livre de tipo finito e um anel de Dedekind. A categoria dos \(L\)-feixes construtíveis é isomorfa à categoria quociente da categoria dos \(A\)-feixes construtíveis pela dos \(A\)-feixes construtíveis de torção. Isso permite, como em (1.4.5), definir, a partir do produto tensorial de (1.4.4), um produto tensorial:
\[
L\text{-}\fc(X)\times L\text{-}\fc(X)\longrightarrow L\text{-}\fc(X).
\]
Como o anel \(A\) é de Dedekind, o argumento usado em (1.4.6) mostra que todo \(A\)-feixe constante torcido é soma direta de um \(A\)-feixe localmente livre e um \(A\)-feixe de torção. Denotando então por \(L\text{-}\fct(X)\) a categoria dos \(L\)-feixes constantes torcidos construtíveis, define-se, como em (1.4.6), um bifuntor
\[
\cHom:(L\text{-}\fct(X))^\circ\times L\text{-}\fc(X)
\longrightarrow L\text{-}\fc(X),
\]
que verifica as propriedades de exatidão habituais.
